\section{FOUR by universal property}

The keystone identifications of this chapter are definitional
(\texttt{Charge} $:=$ \texttt{Articulating.Corner}) or enumerated. This section
records the structure those identifications rely on and the theorem that forces
it: FOUR is initial among bounded interlaced bilattices, so a depth that
exhibits the structure earns its four-corner identification rather than
inheriting it by definition.

\begin{proposition}[Bounded interlaced bilattice]\label{thm:interlaced-bilattice}\lean{Obsmetrologia.InterlacedBilattice}\leanok\uses{thm:four-bilattice}
An \texttt{InterlacedBilattice} is one carrier with two bounded-lattice
structures --- knowledge $(\wedge_k, \vee_k, \bot_k, \top_k)$ and truth
$(\wedge_t, \vee_t, f, t)$ --- each of whose operations is monotone with
respect to the other order. Order is algebraic
($a \le_k b := a \vee_k b = b$, likewise for $\le_t$), and the four
cross-monotonicity fields are stated in that form; own-order monotonicity is
derivable from the lattice laws and is not a field.
\end{proposition}

\begin{theorem}[The bound-exchange identities]\label{thm:bound-exchange}\lean{Obsmetrologia.InterlacedBilattice.joinK_tTop_tBot}\leanok\uses{thm:interlaced-bilattice}
In every bounded interlaced bilattice the bounds of each order are the lattice
combinations of the other order's bounds: $t \vee_k f = \top_k$
(\texttt{joinK\_tTop\_tBot}), $t \wedge_k f = \bot_k$
(\texttt{meetK\_tTop\_tBot}), $\top_k \vee_t \bot_k = t$
(\texttt{joinT\_kTop\_kBot}), and $\top_k \wedge_t \bot_k = f$
(\texttt{meetT\_kTop\_kBot}). Interlacing alone forces these: monotonicity of
the cross-order operations squeezes the combination between the target bound
from both sides. The four bounds therefore form a sub-bilattice.
\end{theorem}

\begin{theorem}[FOUR is initial]\label{thm:four-initial}\lean{Obsmetrologia.four_initial}\leanok\uses{thm:bound-exchange, thm:four-object}
Into every bounded interlaced bilattice $B$ the bounds map
(\texttt{boundsMap}: each charge to the like-named bound) preserves both
lattice structures --- all four binary operations, by
\texttt{boundsMap\_joinK}/\texttt{meetK}/\texttt{joinT}/\texttt{meetT}, whose
sixteen cases each reduce to a bound-exchange identity or a derived
bound law --- and any map out of FOUR sending the four bounds to the four
bounds is the bounds map pointwise (\texttt{boundsMap\_unique}): FOUR's
elements are its bounds, so bound preservation forces the whole map. Existence
and uniqueness together are initiality.
\end{theorem}

\begin{theorem}[Rigidity at four]\label{thm:four-rigid}\lean{Obsmetrologia.four_rigid}\leanok\uses{thm:four-initial}
If the four bounds of $B$ are pairwise distinct, the bounds map is injective
(\texttt{boundsMap\_injective}); if the carrier holds only the four bounds, it
is surjective (\texttt{boundsMap\_surjective}). A bounded interlaced bilattice
with exactly four pairwise-distinct bound elements receives the bounds map as
a bijection: it is FOUR up to the unique bounds-preserving map.
\end{theorem}

\begin{proposition}[FOUR inhabits its own universal property]\label{thm:charge-interlaced}\lean{Obsmetrologia.chargeInterlaced}\leanok\uses{thm:interlaced-bilattice, thm:four-bilattice}
\texttt{chargeInterlaced} equips \texttt{Charge} itself with the
\texttt{InterlacedBilattice} structure, and \texttt{boundsMap\_charge} shows
the unique bounds-preserving map from FOUR to itself is the identity,
pointwise. All declarations of this section are axiom-free.
\end{proposition}

\begin{theorem}[Each depth earns its four]\label{thm:depth-earns-four}\notready\uses{thm:four-initial, thm:articulation-phase-is-charge}
For each depth of the spine that meets the recurring four --- the nuclear
quartet of fixed-point algebras, the bigraded tower, the modal corners, the
regime lattice --- exhibit a bounded interlaced bilattice on that depth's
carrier with pairwise-distinct bounds and derive the depth's four-corner
identification from \ref{thm:four-initial} and \ref{thm:four-rigid}, replacing
the enumerated bijections of the keystone section. A depth whose carrier
provably cannot carry the structure (only one order, or coincident bounds) is
a finding, not a failure: it locates where the recurrence genuinely stops.
\end{theorem}
